\documentclass[a4paper,10pt]{article}
\usepackage[utf8]{inputenc}

\usepackage{amsmath,amssymb,amsfonts,dsfont}
\usepackage{algorithm2e}
\usepackage{tikz}
\usepackage{upgreek}
\bibliographystyle{alpha}
\title{Frobenius Additive Fast Fourier Transform}
\author{}

\begin{document}
\maketitle
\begin{section}{Introduction}
\begin{subsection}{Multiplication 101}
    The multiplication operation is at the core of many scientific applications. The most common way of performing multiplication is to use the
    \textit{classical} method, which involces multiplying the two operands term by term and then adding up the results. This method is easy to understand
    and to apply, however it is very inefficient, especially for large operands: if the operands have $n$ terms, the classical method requires $O(n^2)$
    operations to calculate the product.
    


    There are however other methods to perform multiplication with a better asymptotic complexity, such as the ones based on the Fast Fourier Transform.
    \begin{subsubsection}{Multiplication with Fourier Transforms}
    If we think of the $n$ digits of each input as coefficients of a polynomial, we can calculate the product by following these steps:
    \begin{enumerate}
        \item Sample the polynomials at $2n$ points
        \item Perform the pointwise multiplication of the evaluations
        \item Interpolate the result to obtain the coefficients of the product
    \end{enumerate}

    For example, to calculate $123\cdot 45$:
    \begin{align*}
        a = 123 &\longrightarrow a(x) = 1x^2 + 2x + 3 \\
        b =  45 &\longrightarrow b(x) = 0x^2 + 4x + 5
    \end{align*}
    We sample the polynomials at $2n = 6$ points:
    \begin{align*}
        a(0) = &1\cdot 0^2 + 2\cdot 0 + 3 = 3 \quad &b(0) = &0\cdot 0^2 + 4\cdot 0 + 5 = 5 \\
        a(1) = &1\cdot 1^2 + 2\cdot 1 + 3 = 6 \quad &b(1) = &0\cdot 1^2 + 4\cdot 1 + 5 = 9 \\
        a(2) = &1\cdot 2^2 + 2\cdot 2 + 3 = 11 \quad &b(2) = &0\cdot 2^2 + 4\cdot 2 + 5 = 13 \\
        a(3) = &1\cdot 3^2 + 2\cdot 3 + 3 = 18 \quad &b(3) = &0\cdot 3^2 + 4\cdot 3 + 5 = 17 \\
        a(4) = &1\cdot 4^2 + 2\cdot 4 + 3 = 27 \quad &b(4) = &0\cdot 4^2 + 4\cdot 4 + 5 = 21 \\
        a(5) = &1\cdot 5^2 + 2\cdot 5 + 3 = 38 \quad &b(5) = &0\cdot 5^2 + 4\cdot 5 + 5 = 25 \\
    \end{align*}
    The pointwise multiplication of the evaluations is:
    \begin{align*}
        c(0) = &a(0)\cdot b(0) = 3\cdot 5 = 15 \\
        c(1) = &a(1)\cdot b(1) = 6\cdot 9 = 54 \\
        c(2) = &a(2)\cdot b(2) = 11\cdot 13 = 143 \\
        c(3) = &a(3)\cdot b(3) = 18\cdot 17 = 306 \\
        c(4) = &a(4)\cdot b(4) = 27\cdot 21 = 567 \\
        c(5) = &a(5)\cdot b(5) = 38\cdot 25 = 950 \\
    \end{align*}
    Finally, we can interpolate the result to obtain the coefficients of the product:
    \begin{equation*}
        c(x) = 4x^3 + 13x^2 + 22x + 15
    \end{equation*}
    and convert back to an integer by shifting and adding the coefficients, finding that $123\cdot 45 = 5535$.
\end{subsubsection}


    This approach performs the pointwise multiplication in $\mathcal{O}(n)$, but requires $\mathcal{O}(n^2)$ operations to perform the sampling step
    and $\mathcal{O}(n^3)$ operations to interpolate the result, making it worse than the classical method.

    We are however free to choose the points at which we sample the polynomials, which can drastically reduce the amount of operations needed.
    Instead of sampling at $2n$ random points, we can choose the points to be the $N$-th roots of unity, \textit{i.e.} we calculate the Discrete
    Fourier Transform of the input polynomials:
    \begin{equation}
        \mathcal{F}(x) = \sum_{n=0}^{N} x_n e^{-\frac{i2\pi}{N}kn}
    \end{equation}\label{DFT}
    At this point we can perform the pointwise product in the \textit{transformed} domain and, since $\mathcal{F}$ is invertible, apply the inverse transform:
    \begin{equation}
        \mathcal{F}^{-1}(X) = \frac{1}{N}\sum_{n=0}^{N} X_n e^{i\frac{2\pi}{N}kn}
    \end{equation}
    The algorithms that compute the discrete Fourier Transform of a sequence (or its inverse) are called fast Fourier transforms (FFT), and have an asymptotic complexity
    of $\mathcal{O}(n\log n)$.

\end{subsection}
\begin{subsection}{Segmentation of large binary polynomials}
    Suppose we are multiplying two binary polynomials $f(x)$ and $g(x) \in \mathds{F}_2[x]$, represented as sequences of $d$ bits:
    \begin{equation*}
        a(x) = a_0 + a_1x + \cdots + a_{d-1}x^{d-1}
    \end{equation*}
    In some cases, especially with high degree polynomials, it can be useful to split the polynomials into smaller segments, which can be easier to handle. 
    The standard \textit{Kronecker segmentation} for multiplying binary polynomials is performed as follows
    \begin{enumerate}
        \item Split the polynomials into $n = \lceil d/w \rceil$ $w$-bit blocks
        \begin{align*}
            a(x) &= a_0 + a_1x + \cdots + a_{d-1}x^{d-1} \\
            &= (a_0 + \cdots + a_{w-1}x^{w-1}) + \cdots + (a_{d-w} + \cdots + a_{d-1}x^{w-1})x^{w(n-1)}
        \end{align*}

        \item Define the field $\mathds{F}_{2^{2w}}:=\mathds{F}_2[z]/g(z)$, with $g(z)$ irreducible polynomial of degree $2w$.
        \item Define the map $\psi: \mathds{F}_2[x] \rightarrow \mathds{F}_{2^{2w}[y]}$ such that:
        \begin{equation*}
            a^{\prime}(y) = \psi(a(x)) := a_0^{\prime} + a^{\prime}_1y + \cdots + a^{\prime}_{n-1}y^{n-1}
        \end{equation*}
        where $a^{\prime}_i = (a_{iw} + a_{iw + 1}z + \cdots + a_{iw+(w-1)}z^{w-1})$.
    
        \item Calculate the product $c^{\prime}(y) = a^{\prime}(y)\cdot b^{\prime}(y) = c_0^{\prime} + c^{\prime}_1y + \cdots \in \mathds{F}_{2^{2w}}[y]$ \label{list:mult}
        \item Map $z \mapsto x$ and $y \mapsto x^w$, collect terms and coefficients and find the result.
    \end{enumerate}

    We know that it is possible to perform the multiplication at step~\eqref{list:mult} with FFT.
    Following the approach in~\cite{thomas2016introduction}, the steps are:
    \begin{enumerate}
        \item (\texttt{fft}) Evaluate the polynomials $a^{\prime}(y)$ and $b^{\prime}(y)$ at $2n$ points using FFT.
        \item (\texttt{pointmul}) Multiply the evaluated values pairwise.
        \item (\texttt{ifft}) Interpolate back into a polynomial of degree at most $2n-1$ using inverse FFT.
    \end{enumerate}
\end{subsection}
\end{section}


\section*{References}
\bibliography{refs}

\end{document}

